% glossary.tex -- Module 6 ESSENTIALS glossary (A-Z). GENERATED by
% src/mod6/glossary.tex (the key terms only; meanings and Chinese unchanged,
% a few meanings shortened). Edit the script, not this file.
\unnumberedchapter[\faBook]{Glossary}

The key words of this module, from A to Z. The meaning is in plain English. The last
column gives the usual Chinese term used in physics textbooks. The full guide has a
longer glossary.

{\small
\begin{glossarytable}
  % ---- A
  \gterm{AC induction motor}{a motor whose rotor currents are induced by a rotating field}{交流感应电动机（异步电动机）}
  \gterm{Alternating current (AC)}{current that keeps reversing its direction}{交流电}
  \gterm{Ampere (A)}{the SI unit of current: \qty{1}{A} means one coulomb of charge passes a point each second}{安培}
  % ---- B
  \gterm{Back emf ($\varepsilon_\mathrm{back}$)}{the emf induced in a spinning motor coil; it opposes the supply}{反电动势}
  \gterm{Brushes}{fixed contacts that join the spinning coil to the circuit}{电刷}
  % ---- C
  \gterm{Centripetal force}{the net force towards the centre of a circle}{向心力}
  \gterm{Charge-to-mass ratio ($q/m$)}{charge divided by mass; a bigger $q/m$ gives a bigger acceleration}{荷质比（比荷）}
  \gterm{Conventional current}{the way positive charge would flow: from $+$ round the circuit to $-$}{（规定的）电流方向}
  \gterm{Core (soft-iron core)}{the iron that carries the magnetic field from one coil to the other}{铁芯（软铁芯）}
  \gterm{Couple}{two equal and opposite forces that turn an object without moving it sideways}{力偶}
  % ---- D
  \gterm{DC motor}{a device that turns electrical energy from a DC supply into rotational kinetic energy}{直流电动机}
  \gterm{Deflection}{how far a particle is pushed sideways from its straight path}{偏转（偏移量）}
  \gterm{Direct current (DC)}{current that flows in one direction only}{直流电}
  % ---- E
  \gterm{Eddy currents}{loops of current induced inside solid metal when the flux through it changes}{涡流}
  \gterm{Effective value (of AC)}{the steady (DC) voltage that would heat a heater by the same amount}{有效值}
  \gterm{Efficiency ($\eta$)}{the percentage of the input power that comes out as useful power}{效率}
  \gterm{Electric field}{a region where a charge feels an electric force}{电场}
  \gterm{Electric field strength ($E$)}{the force on each coulomb of charge; unit \si{V.m^{-1}}, the same as \si{N.C^{-1}}}{电场强度}
  \gterm{Electric potential energy ($U$)}{energy a charge stores because of its position in an electric field}{电势能}
  \gterm{Electromagnetic induction}{making an emf by changing the magnetic flux through a circuit}{电磁感应}
  \gterm{Electromotive force (emf, $\varepsilon$)}{not a force: the energy a source gives each coulomb of charge; unit volt (V)}{电动势}
  \gterm{Electron volt (eV)}{the energy one electron gains through \qty{1}{V}: \qty{1}{eV} $=$ \qty{1.602e-19}{J}}{电子伏特}
  \gterm{Elementary charge ($e$)}{the size of the charge on one proton or electron, \qty{1.602e-19}{C}}{元电荷}
  % ---- F
  \gterm{Faraday's law}{the induced emf equals the number of turns times the rate of change of flux}{法拉第电磁感应定律}
  % ---- G
  \gterm{Generator}{a device that turns kinetic energy into electrical energy}{发电机}
  % ---- I
  \gterm{Ideal transformer}{a perfect transformer that wastes no energy, so power in equals power out}{理想变压器}
  \gterm{Incomplete flux linkage (flux leakage)}{some of the flux from the primary misses the secondary}{漏磁（磁通未完全耦合）}
  \gterm{Induced current}{the current that an induced emf drives around a closed circuit}{感应电流}
  \gterm{Induced emf}{an emf produced by a changing magnetic flux}{感应电动势}
  % ---- L
  \gterm{Laminated core}{a core made of thin iron sheets with insulation between them}{叠片铁芯（硅钢片铁芯）}
  \gterm{Lenz's law}{an induced current's magnetic field opposes the change in flux that made it}{楞次定律}
  \gterm{Lever arm}{the perpendicular distance from the axis to the line of a force}{力臂}
  % ---- M
  \gterm{Magnetic braking}{slowing moving metal with eddy currents, without touching it}{电磁制动（电磁阻尼）}
  \gterm{Magnetic field}{a region where a magnet or a moving charge feels a magnetic force}{磁场}
  \gterm{Magnetic field strength ($B$)}{how strong a magnetic field is; also called magnetic flux density; unit tesla (T)}{磁感应强度}
  \gterm{Magnetic flux ($\Phi$)}{how much magnetic field passes through an area; unit weber (Wb)}{磁通量}
  \gterm{Magnetic flux density}{another name for $B$: the flux through each square metre at right angles to the field}{磁通密度（即磁感应强度）}
  \gterm{Magnetic permeability constant ($\mu_0$)}{the constant that links a current to the field it makes in a vacuum, $4\pi \times 10^{-7}$ \si{N.A^{-2}}}{真空磁导率}
  \gterm{Mass spectrometer}{a machine that sorts ions by mass, using the radius of their path in a magnetic field}{质谱仪}
  \gterm{Motional emf}{the emf in a conductor because it moves across a magnetic field}{动生电动势}
  \gterm{Motor effect}{the force on a current-carrying wire in an external magnetic field}{磁场对电流的作用（安培力）}
  \gterm{Mutual induction}{an emf induced in one coil by a changing current in a nearby coil}{互感}
  % ---- N
  \gterm{Normal}{an imaginary line at \ang{90} to a flat surface, such as the plane of a coil}{法线}
  % ---- P
  \gterm{Parabola}{the curved path of a thrown ball, made by a constant force in one direction}{抛物线}
  \gterm{Perpendicular}{at right angles (\ang{90})}{垂直}
  \gterm{Primary coil}{the input coil, joined to the supply}{原线圈（初级线圈）}
  % ---- R
  \gterm{Radial magnetic field}{a field whose lines point along radii, like the spokes of a wheel}{辐向磁场}
  \gterm{Regenerative braking}{braking in which the motor acts as a generator and the battery stores the energy}{再生制动（能量回收制动）}
  \gterm{Resistive heating}{heat made when a current flows through a resistance ($I^2R$)}{电阻发热（焦耳热）}
  \gterm{Right-hand grip rule}{thumb along the current; the curled fingers show the field around it}{右手螺旋定则（安培定则）}
  \gterm{Right-hand palm rule}{fingers along $B$, thumb along $I$ (or $v$ of a positive charge): the palm pushes along $F$}{右手掌定则（结果同左手定则）}
  \gterm{Rotating magnetic field}{a field of fixed strength whose direction turns steadily, made by coils that do not move}{旋转磁场}
  \gterm{Rotor (armature)}{the part of a motor that turns: a coil on an iron core, fixed to the axle}{转子（电枢）}
  % ---- S
  \gterm{Secondary coil}{the output coil, in which the emf is induced}{副线圈（次级线圈）}
  \gterm{Slip}{the difference between the speed of the rotating field and the speed of the rotor}{转差}
  \gterm{Slip rings}{two full rings, each always joined to the same end of a generator coil}{滑环}
  \gterm{Solenoid}{a long coil of wire with many turns; the field inside is uniform}{螺线管}
  \gterm{Split-ring commutator}{two half-rings that reverse the current in the coil every half turn}{换向器（裂环换向器）}
  \gterm{Squirrel-cage rotor}{metal bars joined at both ends by end rings, like a cage}{鼠笼式转子}
  \gterm{Stator}{the part of a motor that does not move; it makes the magnetic field}{定子}
  \gterm{Step-down transformer}{a transformer with fewer secondary turns, so it makes the voltage smaller}{降压变压器}
  \gterm{Step-up transformer}{a transformer with more secondary turns, so it makes the voltage bigger}{升压变压器}
  \gterm{Substation}{a place with transformers that change the voltage of the grid}{变电站}
  % ---- T
  \gterm{Tesla (T)}{the SI unit of magnetic field strength: \qty{1}{T} $=$ \qty{1}{N.A^{-1}.m^{-1}} $=$ \qty{1}{Wb.m^{-2}}}{特斯拉}
  \gterm{Three-phase AC}{three AC supplies that peak one after another, \ang{120} apart}{三相交流电}
  \gterm{Torque ($\tau$)}{the turning effect of a force; unit \si{N.m}}{力矩}
  \gterm{Transformer}{a device that changes an AC voltage}{变压器}
  \gterm{Transmission line}{a long, thick wire on tall towers that carries electricity a long way}{输电线}
  \gterm{Turns ratio}{$N_\mathrm{p} : N_\mathrm{s}$, the primary turns compared with the secondary turns}{匝数比}
  % ---- U
  \gterm{Uniform field}{a field with the same strength and direction everywhere}{匀强场（匀强电场、匀强磁场）}
  % ---- V
  \gterm{Velocity selector}{crossed fields that let only one speed pass straight through}{速度选择器}
  \gterm{Voltage drop}{the voltage used up pushing a current through a resistance, $IR$}{电压降（电压损失）}
  % ---- W
  \gterm{Weber (Wb)}{the SI unit of magnetic flux: \qty{1}{Wb} $=$ \qty{1}{T.m^2}}{韦伯}
\end{glossarytable}}
